%% =========================================================================
%% Data section (2026-07-12): promoted from the former Methods subsection
%% "Data acquisition and preparation" and placed before the Demolition
%% indicators section, which relies on the cleaned extract it describes.
%% Retains \label{sec:data} so existing cross-references resolve.
%% =========================================================================
\section{Data acquisition and preparation}\label{sec:data}

The analysis is based on a complete extract of the BBR building data
and associated case objects, including all historical versions and
lifecycle statuses (Section~\ref{sec:bbr-structure}), obtained from the
Danish Data Distributor (Datafordeler) on 11~July~2026. A Python/Polars
pipeline, archived with the repository
(Section~\ref{sec:data-availability}), parses the raw extract, derives
per-building version histories, classifies building-use codes as
current or discontinued~\cite{bbr_kodelister}, and links lifecycle-exit
transitions to demolition cases.

The study uses three date windows, each for a different purpose. The
\emph{full reporting window} is 2000--2025. Its start is set by the
register itself: the earliest effect date in the building data is
5~February~2000, so no register-dated building event precedes the
window. A thin tail of case dates reaches further back and is excluded.
The end of the window restricts the analysis to complete calendar years:
events dated in 2026 are excluded. This full window is used to describe
the complete set of candidate demolitions observable in the extract and
to show the long annual series.

The \emph{main comparison window} is 2018--2025. These are the first
complete calendar years after the 2~June~2017 Datafordeler registration
floor, so they are used for averaged annual comparisons, windowed
indicator overlap, and the headline demolition-rate comparison. Earlier
years are still shown, but they are interpreted as historical context
because they are affected by backdated effect dates, the BBR system
migration, and a registration step change around 2011. Finally, the
\emph{KMD comparison window} is 2011--2019, used only when comparing to
the Andersen/KMD extract: the source study discards pre-2011 records and
the extract ends during 2020, making 2020 a partial year.

\subsection{Field completeness and data quality}\label{sec:missingness}
%% Moved here 2026-07-12 from Results (former sec:missingness). Descriptive
%% characterization of the extract; the demolished-area SPREAD it sets up is
%% the result, reported in Section~\ref{sec:area-sensitivity}. Numbers from
%% docs/data_quality_report.md / results/data_quality_*.csv (regenerated by
%% `python -m src.data_quality`). Scope decision 2026-07-12: report byg038 +
%% byg041 only, on ALL currently effective buildings (no status-10 slice).

The building register comprises 33{,}820{,}994 temporal rows for
6{,}250{,}075 buildings, a mean of 5.41 versions per building; the linked
case and case-level objects add 6{,}507{,}234 and 8{,}426{,}674 rows
respectively. All three share a registration floor of 2~June~2017, the
date of the BBR~1.8 system migration, so events before that date are
represented only through backdated effect dates
(Section~\ref{sec:bbr-structure}).

The fields the indicators and area measures depend on are largely
complete: the business-process code (\texttt{forretningsproces}) is
missing for 0.02\,\% of building rows and the building-use code
(\texttt{byg021}) for 0.49\,\%. The construction year is missing for
15.81\,\%, but it is carried only as a covariate and is never aggregated,
so its gaps do not affect the demolition and area outcomes.

Table~\ref{tab:data-quality-constraints} summarizes the quality checks
that affect the demolition indicators directly. The purpose is not to
audit every BBR field, but to show which limitations can change the
candidate demolition set or its dating. The building-level signals are
not limited by missing field coverage: both \texttt{forretningsproces}
and \texttt{sagstype} are almost complete. The main limitation for the
case-based indicators is instead linkage. Among the 599{,}431 temporal
case-level rows marked as partial or total demolition
(\texttt{sagstype} 31 or 32), 82{,}214 rows (13.7\,\%) lack a
\texttt{stamdataBygning} reference and therefore cannot be attached to a
building. The linkable demolition-case rows identify 257{,}602 unique
buildings. Nearly all of these buildings have a case date
(\texttt{sag002 Byggesagsdato}; 97.1\,\%), but the case completion field
(\texttt{sag010 Fuldf{\o}relseAfByggeri}) is not used as a demolition
completion date. Although it is present for 84.6\,\% of linkable
demolition-case buildings, it belongs to the joined building-case record
and can date co-filed construction or alteration work rather than the
demolition itself (Section~\ref{sec:bbr-chain}).

\begin{table}[htbp]
\centering
\small
\caption{Analysis-relevant data-quality constraints in the cleaned BBR
extract. Field-level missingness is measured on currently effective rows;
demolition-case linkage and dates are measured on temporal
\texttt{Sagsniveau} rows with \texttt{sagstype} 31 or 32.}
\label{tab:data-quality-constraints}
\begin{tabular}{@{}p{0.26\linewidth}p{0.31\linewidth}p{0.34\linewidth}@{}}
\toprule
Constraint & Observed in extract & Consequence for the analysis \\
\midrule
Registration-history floor &
All three entities share the BBR~1.8 registration floor,
2~June~2017. &
Pre-2017 events are observed through backdated effect dates and are
therefore shown separately from the main 2018--2025 comparison window. \\
\addlinespace
Building-level process code &
\texttt{forretningsproces} missing for 0.02\,\% of building rows. &
D2 is not mainly limited by missing process-code coverage; disagreement
with other indicators reflects the signal, not field absence. \\
\addlinespace
Case type &
\texttt{sagstype} missing for 0.08\,\% of case-level rows. &
The demolition-case family is not materially limited by missing
case-type coverage. \\
\addlinespace
Demolition-case linkage &
599{,}431 temporal rows have \texttt{sagstype} 31 or 32; 82{,}214
(13.7\,\%) lack \texttt{stamdataBygning}. &
Case-based indicators D4--D6 can count only demolition cases linkable to
a building. \\
\addlinespace
Case dating &
Of 257{,}602 linkable demolition-case buildings, 250{,}169 (97.1\,\%)
have \texttt{sag002}; 218{,}000 (84.6\,\%) have \texttt{sag010}. &
\texttt{sag002} can date case activity; \texttt{sag010} is not treated
as a demolition-completion date because it is not demolition-specific. \\
\addlinespace
Total floor area &
\texttt{byg038} is missing for 51.24\,\% of currently effective
buildings and for 54.84\,\% of status-10 buildings. &
Demolished-area totals must be reported with coverage; missingness is
not random across building uses. \\
\bottomrule
\end{tabular}
\end{table}

The floor-area measures are the central data-quality constraint on the
area axis (Section~\ref{sec:areas}). Across all currently effective
buildings, the footprint \texttt{byg041} is missing for 9.29\,\%, but
the total building area \texttt{byg038} is missing for 51.24\,\%. The
only near-complete area column is therefore a \emph{different physical
quantity}---ground footprint, a single storey---from the total floor
area that the comparison literature reports.

This missingness is structural rather than random. BBR does not permit a
total floor area for the smallest building-use codes (910--930), so
\texttt{byg038} is missing for 99.78\,\% of outbuildings and other minor
structures, against about 11\,\% for housing. The aggregate 51.24\,\% is
dominated by these small structures: among the 3.5~million currently
effective buildings with a known use code \emph{outside} the outbuilding
group, 13.3\,\% lack a total floor area (a further 2.2\,\% record a
value of zero). Any use of \texttt{byg038} therefore covers a
use-dependent, non-random subset of the stock---a coverage effect that
the total-floor-area coverage analysis quantifies
(Section~\ref{sec:area-sensitivity}).
